\subsection{A concrete obstruction to the finite-Frattini extension}

\begin{claim}{A.R6}{Proved}{Mixed-prime surface groups have infinite Frattini quotient}
Let $\Sigma$ contain two distinct primes $r,s$, and let $\Delta$
be the pro-$\Sigma$ completion of a surface group of genus at least
two. Its Frattini quotient $\Delta/\Phi(\Delta)$ is infinite,
where $\Phi$ is the intersection of the maximal open subgroups.
In particular this applies to the geometric pro-$\Sigma$ group
of a proper curve in characteristic outside $\Sigma$.

For $\Sigma=\{2,3\}$ there is, moreover, an explicit two-generated
group $G$ of order $576$, with $\Phi(G)=1$, and an automorphism of
order seven acting trivially on $G^{\mathrm{ab}}$.
Neither assertion is a counterexample to mixed-prime point separation.
They show why the finite-Frattini and abelianization arguments used
for a single prime do not extend by the same reasoning.
\end{claim}
\begin{proof}
For each positive integer $n$, put $d_n=\operatorname{ord}_{s^n}(r)$.
Choose an element $\zeta_n$ of exact order $s^n$ in
$\mathbf F_{r^{d_n}}^\times$, and form
\[
 G_n=(\mathbf F_{r^{d_n}},+)\rtimes\langle\zeta_n\rangle,
 \qquad (a,i)(b,j)=(a+\zeta_n^i b,i+j).
\]
The action on the additive normal subgroup is irreducible over
$\mathbf F_r$, because $\mathbf F_r[\zeta_n]=\mathbf F_{r^{d_n}}$.
It is faithful. Consequently the complement $H_n=\langle\zeta_n\rangle$
is a maximal subgroup: a subgroup containing it intersects the
additive group in an invariant subspace. Its core is trivial,
since a nonidentity scalar fails to fix some additive element.
Thus $\Phi(G_n)\subset\bigcap_{g\in G_n}gH_ng^{-1}=1$.
The elements $(1,0)$ and $(0,1)$ generate $G_n$, since the scalar
orbit of $1$ spans the whole additive group.

The surface group surjects onto the free group on two generators
by sending two of its $a$-generators to free generators and every
$b$-generator, and all remaining $a$-generators, to one. Its
pro-$\Sigma$ completion therefore surjects onto each $G_n$.
The image of a Frattini subgroup under a surjection is contained
in the target Frattini subgroup: pull back every maximal subgroup.
Hence $\Delta/\Phi(\Delta)$ surjects onto every $G_n$.
Their orders are unbounded, since they are divisible by $s^n$;
the quotient is infinite. The characteristic-$p$ interpretation
uses the geometric specialization isomorphism already checked in
A.CF3, with all primes of $\Sigma$ invertible.

For the finite example use
\[
 \mathbf F_{64}=\mathbf F_2[t]/(t^6+t+1),\quad
 \zeta=t^7,\quad\lambda=t^9,
 \qquad G=(\mathbf F_{64},+)\rtimes\langle\zeta\rangle.
\]
Here $\zeta$ has order nine and generates the degree-six field;
$\lambda$ has order seven. Multiplication by $\lambda$ on the
additive factor commutes with multiplication by $\zeta$, so
\[
                    (a,i)\longmapsto(\lambda a,i)
\]
is an automorphism of order seven. Since $\zeta-1$ is invertible,
the commutators with $\zeta$ fill the additive factor. Thus
$G^{\mathrm{ab}}=\mathbf Z/9$, on which this automorphism is
the identity. The program \path{code/frattini_two_prime.py}
checks the field, the root orders, all $576^2$ automorphism
compatibilities, the two-generator orbit, the commutator image,
and the trivial complement core. This supplies the exact finite
counterexample to replacing the Frattini quotient by abelianization
in the attempted extension of the single-prime argument.
\end{proof}
